% !PRINTED 89
\weberpage{89}
%<<CONT

7. \gesperrt{Es muss die Anzahl} $\varrho$ \gesperrt{der durch eine Primzahl} $p$
\gesperrt{theilbaren Invarianten gleich oder kleiner als die Anzahl} $\mu$
\gesperrt{der Indexmoduln sein, und es muss sich aus der Matrix}

\weq{24}{
\begin{gathered}
e_{1,1}\,i_{1,1}, \quad e_{2,1}\,i_{2,1}, \; \dots \; e_{\varrho,1}\,i_{\varrho,1} \\
e_{1,2}\,i_{1,2}, \quad e_{2,2}\,i_{2,2}, \; \dots \; e_{\varrho,2}\,i_{\varrho,2} \\
\cdot\ \cdot\ \cdot\ \cdot\ \cdot\ \cdot\ \cdot\ \cdot\ \cdot\ \cdot \\
e_{1,\mu}\,i_{1,\mu}, \quad e_{2,\mu}\,i_{2,\mu}, \; \dots \; e_{\varrho,\mu}\,i_{\varrho,\mu}
\end{gathered}
}

\gesperrt{wenigstens eine durch} $p$ \gesperrt{untheilbare Determinante von}
$\varrho$ \gesperrt{Reihen bilden lassen.}

Dies involvirt zun\"achst eine Forderung f\"ur die $i_{k,h}$, d.\,h. also
in letzter Instanz eine Bedingung f\"ur die Zahl $m$.

Alle Glieder irgend einer $\varrho$-reihigen Determinante der Matrix (24)
enthalten einen Factor der Form

\[
i_{1,h_1}\,i_{2,h_2}, \; \dots \; i_{\varrho,h_\varrho},
\]

worin $h_1, h_2, \dots h_\varrho$ irgend eine Combination von $\varrho$
verschiedenen Zahlen der Reihe $1, 2, \dots \mu$ bedeuten. W\"aren nun alle
diese Producte durch $p$ theilbar, so w\"aren sicher alle Determinanten der
Matrix (24) von $\varrho$ Reihen auch durch $p$ theilbar, und es muss also
wenigstens eine solche Combination geben, die nicht durch $p$ theilbar ist.
Weil aber die $i_{k,h}$ nur Potenzen von $p$ sein k\"onnen, so muss

\[
i_{1,h_1} = 1, \quad i_{2,h_2} = 1, \; \dots \; i_{\varrho,h_\varrho} = 1
\]

sein. Geht man nun auf die Bedeutung der $i_{k,h}$ in (11) zur\"uck, so
k\"onnen wir die letzte Forderung f\"ur den Grad $m$ folgendermaassen
ausdr\"ucken:

f) \gesperrt{Wenn eine Primzahl} $p$ \gesperrt{in} $\varrho$
\gesperrt{Invarianten} $i_1, i_2, \dots i_\varrho$ \gesperrt{aufgeht, so muss
die Anzahl} $\mu$ \gesperrt{der Indexmoduln gleich oder gr\"osser als}
$\varrho$ \gesperrt{sein, und es m\"ussen sich} $\varrho$ \gesperrt{dieser
Indexmoduln} $c_{h_1}, c_{h_2}, \dots c_{h_\varrho}$ \gesperrt{so
ausw\"ahlen lassen, dass} $c_{h_1}$ \gesperrt{durch} $i_1$, $c_{h_2}$
\gesperrt{durch} $i_2 \dots$, $c_{h_\varrho}$ \gesperrt{durch} $i_\varrho$
\gesperrt{theilbar ist} $^1$).

\noindent\rule{4cm}{0.4pt}

\footnotesize
$^1$) Der Nachweis der Thatsache, dass dieser Forderung immer auf
unendlich viele Arten entsprochen werden kann, st\"utzt sich auf den Satz
der Zahlentheorie, dass unter den Zahlen einer arithmetischen Progression,
deren Anfangsglied und Differenz ohne gemeinsamen Theiler sind, unendlich
viele Primzahlen vorkommen, ein Satz, der bis jetzt nur von
\normalsize
%CONT>>
